
\documentclass{article}
\usepackage{colortbl}
\usepackage{makecell}
\usepackage{multirow}
\usepackage{supertabular}

\begin{document}

\newcounter{utterance}

\centering \large Interaction Transcript for game `referencegame', experiment `letter\_grids', episode 11 with reflexion{-}Qwen3.5{-}9B.
\vspace{24pt}

{ \footnotesize  \setcounter{utterance}{1}
\setlength{\tabcolsep}{0pt}
\begin{supertabular}{c@{$\;$}|p{.15\linewidth}@{}p{.15\linewidth}p{.15\linewidth}p{.15\linewidth}p{.15\linewidth}p{.15\linewidth}}
    \# & $\;$A & \multicolumn{4}{c}{Game Master} & $\;\:$B\\
    \hline

    \theutterance \stepcounter{utterance}  
    & & \multicolumn{4}{p{0.6\linewidth}}{
        \cellcolor[rgb]{0.9,0.9,0.9}{
            \makecell[{{p{\linewidth}}}]{
                \texttt{\tiny{[P1$\langle$GM]}}
                \texttt{You have attempted this same task before and did not succeed. These are the plans you wrote after those attempts.} \\
\texttt{Trial \#1: New plan: My previous description failed because it incorrectly identified the "bottom{-}left column" as completely filled (it is only filled in the first three rows) and mischaracterized the empty area as a "2x4 area" (the empty region is actually a 3x4 block spanning rows 2{-}4 and columns 2{-}5). To succeed, I must strictly adhere to the visual evidence: describe the top row as fully filled, the leftmost column as filled only for the first three rows, and the third row as fully filled. I will explicitly state that the remaining cells (rows 2{-}4, columns 2{-}5) are empty to ensure the description uniquely distinguishes the target from the distractors, which have different patterns of filled rows and columns.} \\
\\ 
\texttt{Trial \#2: New plan: My previous description failed because it incorrectly stated that the leftmost column is filled only in the first three rows; in reality, the leftmost column is completely filled (rows 1{-}5). Additionally, I misidentified the empty region as spanning rows 2{-}4, when it actually spans rows 2{-}5. To succeed, I must strictly adhere to the visual evidence: describe the top row (row 1) and the third row (row 3) as completely filled, the leftmost column as completely filled, and explicitly state that the remaining cells (rows 2, 4, and 5, columns 2 through 5) are empty. This precise distinction will differentiate the target from Distractor 1 (which has a filled second row) and Distractor 2 (which has a filled fifth row).} \\
\\ 
\\ 
\texttt{You are given three grids, where each of them is 5 by 5 in size.} \\
\texttt{Grids have empty cells marked with "▢" and filled cells marked with "X".} \\
\texttt{Your task is to generate a referring expression that best describes the target grid while distinguishing it from the two other distractor grids.} \\
\texttt{The first grid is the target grid, and the following two grids are the distractors.} \\
\\ 
\texttt{Target grid:} \\
\\ 
\texttt{X X X X X} \\
\texttt{□ □ □ □ □} \\
\texttt{X X X X □} \\
\texttt{X □ □ □ □} \\
\texttt{X □ □ □ □} \\
\\ 
\texttt{Distractor grid 1:} \\
\\ 
\texttt{X X X X X} \\
\texttt{X □ □ □ □} \\
\texttt{X X X X X} \\
\texttt{X □ □ □ □} \\
\texttt{X X X X X} \\
\\ 
\texttt{Distractor grid 2:} \\
\\ 
\texttt{X X X X X} \\
\texttt{X □ □ □ □} \\
\texttt{X □ □ □ □} \\
\texttt{X □ □ □ □} \\
\texttt{X X X X X} \\
\\ 
\texttt{Instruction: Describe the target grid.} \\
\texttt{Generate the referring expression starting with the tag "Expression: " for the given target grid. Omit any other text.} \\
            }
        }
    }
    & \\ \\

    \theutterance \stepcounter{utterance}  
    & \multicolumn{4}{p{0.6\linewidth}}{
        \cellcolor[rgb]{0.9,0.9,0.9}{
            \makecell[{{p{\linewidth}}}]{
                \texttt{\tiny{[P1$\rangle$GM]}}
                \texttt{Expression: The target grid has a completely filled top row (row 1) and a completely filled third row (row 3). Additionally, the leftmost column (column 1) is completely filled from top to bottom. All other cells in the grid are empty. This distinguishes it from Distractor 1, which has filled rows 1, 3, and 5, and Distractor 2, which has filled rows 1, 3, and 5 but with a different pattern in the middle rows (specifically, Distractor 2 has empty cells in column 1 for rows 2, 3, and 4, whereas the target has filled cells in column 1 for all rows). Wait, let me re{-}examine the grids carefully.} \\
\\ 
\texttt{Target:} \\
\texttt{Row 1: X X X X X (Full)} \\
\texttt{Row 2: □ □ □ □ □ (Empty)} \\
\texttt{Row 3: X X X X □ (Not full, last cell empty) {-}> Wait, looking at the text representation: "X X X X □". That means row 3 is NOT fully filled.} \\
\texttt{Row 4: X □ □ □ □ (Only first cell filled)} \\
\texttt{Row 5: X □ □ □ □ (Only first cell filled)} \\
\\ 
\texttt{Let's re{-}read the Target grid provided in the prompt:} \\
\texttt{Target grid:} \\
\texttt{X X X X X} \\
\texttt{□ □ □ □ □} \\
\texttt{X X X X □} \\
\texttt{X □ □ □ □} \\
\texttt{X □ □ □ □} \\
\\ 
\texttt{Row 1: All X.} \\
\texttt{Row 2: All □.} \\
\texttt{Row 3: X X X X □ (First 4 are X, last is □).} \\
\texttt{Row 4: X □ □ □ □ (First is X, rest □).} \\
\texttt{Row 5: X □ □ □ □ (First is X, rest □).} \\
\\ 
\texttt{Distractor 1:} \\
\texttt{X X X X X} \\
\texttt{X □ □ □ □} \\
\texttt{X X X X X} \\
\texttt{X □ □ □ □} \\
\texttt{X X X X X} \\
\\ 
\texttt{Distractor 2:} \\
\texttt{X X X X X} \\
\texttt{X □ □ □ □} \\
\texttt{X □ □ □ □} \\
\texttt{X □ □ □ □} \\
\texttt{X X X X X} \\
\\ 
\texttt{Okay, my previous analysis in the "thought" block was flawed because I misread the target grid's row 3. Let's correct the observation.} \\
\\ 
\texttt{Target Grid Analysis:} \\
\texttt{{-} Row 1: Fully filled (5 Xs).} \\
\texttt{{-} Row 2: Fully empty (5 □s).} \\
\texttt{{-} Row 3: First 4 cells filled, last cell empty.} \\
\texttt{{-} Row 4: Only first cell filled.} \\
\texttt{{-} Row 5: Only first cell filled.} \\
\texttt{{-} Column 1: Fully filled (Rows 1{-}5 are X).} \\
\texttt{{-} Column 2: Row 1 is X, Row 2 is □, Row 3 is X,} \\
            }
        }
    }
    & & \\ \\

    \theutterance \stepcounter{utterance}  
    & & & \multicolumn{2}{p{0.3\linewidth}}{
        \cellcolor[rgb]{0.9,0.9,0.9}{
            \makecell[{{p{\linewidth}}}]{
                \texttt{\tiny{[GM$|$GM]}}
                \texttt{Answer followed by extra content} \\
            }
        }
    }
    & & \\ \\

    \theutterance \stepcounter{utterance}  
    & \multicolumn{4}{p{0.6\linewidth}}{
        \cellcolor[rgb]{0.95,0.95,0.95}{
            \makecell[{{p{\linewidth}}}]{
                \texttt{\tiny{[P2$\rightarrow$P1]}}
                \texttt{Post{-}episode Reflexion (not sent to GM):} \\
\texttt{New plan: My previous attempts failed because I generated verbose explanations and internal monologues instead of providing only the required "Expression: " tag followed immediately by the description. I also made factual errors regarding the grid contents, specifically misidentifying the fill status of Row 3 and the extent of the empty region. To succeed, I will strictly output only the tag "Expression: " followed by a concise, accurate description: "The target grid has a completely filled top row (row 1) and a completely filled third row (row 3), with the leftmost column (column 1) fully filled from top to bottom. All other cells are empty." This description correctly distinguishes the target from Distractor 1 (which has a filled row 2) and Distractor 2 (which has a filled row 5).} \\
            }
        }
    }
    & & \\ \\

\end{supertabular}
}

\end{document}
